% =====================================================================
% related_work.tex — Related Work section, R4 paper draft v0.1 (2026-10-04)
% Distilled from docs/research_survey_v1_en.md (authoritative English edition); all \cite keys match latex/refs.bib
% and refs/manifest.md, and every entry has a verified PDF (refs/pdfs/).
%
% Journal strategy (survey §7):
%   - Minds and Machines: keep the conceptual-analysis detail of §RW.1/§RW.2/§RW.6 (philosophical controversies),
%     compress the engineering-technical description of §RW.5;
%   - Synthese: keep the measurement-methodology detail of §RW.5/§RW.9 (representation theorem, identifiability, dual track),
%     compress the philosophy-history exposition of §RW.1.
%   Scalable passages are marked in both cases by LaTeX comments [MaM] / [Syn].
%
% Build dependency: Springer Nature sn-jnl template (\documentclass[...]{sn-jnl}),
% see latex/README.md.
% =====================================================================

\section{Related Work}\label{sec:related-work}

Our framework sits at the intersection of seven literatures that have
hitherto
developed largely in isolation: philosophical theories of autonomy, cybernetic
and enactive accounts of agency, the empirical psychology of self-regulation,
information-theoretic measures of control, the emerging science of AI agent
characterization and governance, AI safety research on goal corruption, and
the formal study of delegation and adversarial allocation (assistance games,
contract theory, security games).
This section surveys each, identifies the points of consensus and contention
that structure the field (\S\ref{sec:related-work}.1--.8), and states the gaps
our contribution fills (\S\ref{sec:related-work}.9). Throughout, we use the
five-channel vocabulary of Section~\ref{sec:framework} only informally; the
formal definitions are self-contained.

\subsection{Philosophical accounts of autonomy}\label{rw:philosophy}

% [MaM: keep full passage; Syn: may compress to first two + last sentence]
Modern analysis begins with hierarchical accounts. \citet{frankfurt1971}
grounds autonomy in the reflective structure of the will: an agent acts
autonomously when the desire that moves her is endorsed by a second-order
volition with which she identifies. \citet{watson1975} objects that a purely
hierarchical picture invites a regress---nothing guarantees that the endorsing
faculty is itself autonomous---and proposes instead that autonomy consists in
the governance of motivation by the agent's valuational system. Both accounts
are \emph{procedural}: they constrain how motives are formed and reviewed, not
what they contain. Proceduralism was refined by \citet{christman2004} into
conditions of ``potential reflective acceptability'' that avoid demanding
occurrent introspection.

The decisive challenge to proceduralism is the manipulation case: an agent
whose entire motivational structure, endorsement faculty included, was
installed by a designer passes every procedural test while intuitively lacking
autonomy. \citet{fischer_ravizza1998} respond with a \emph{historical}
condition---the mechanism issuing in action must be moderately
reasons-responsive \emph{and} the agent must have taken responsibility for it,
so that provenance, not merely present structure, matters. Whether history is
genuinely required remains contested (Section~\ref{sec:objections}); we return
to it in \S\ref{sec:axioms}, where it becomes Axiom~AX3 and, crucially, an
operationally testable condition (Theorem~T3).

A third strand, relational autonomy, rejects the atomistic picture shared by
both sides. \citet{mackenzie_stoljar2000} argue that the capacities
constitutive of autonomy---self-trust, self-esteem, the standing to avow
commitments---are socially scaffolded; \citet{westlund2009} adds that holding
oneself answerable to others is constitutive rather than derivative. On this
view relationships are the substrate of autonomy, not its principal threat.
Relational autonomy licenses the distinction, central to our intervention
taxonomy, between \emph{constitutive support} and \emph{reductive
interference}: the same causal coupling can raise or lower autonomy depending
on its functional role. Externalism supplies the final coordinate:
\citet{oshana1998} argues that autonomy is partly constituted by
socio-relational and independence conditions external to the agent's
psychology, against purely internalist procedural accounts. Our framework
assigns each camp its instruments rather than adjudicating between them: the
internalist clauses live in the origination and endorsability respects of
Definition~\ref{def:autonomy}, while the externalist conditions are carried
by the viability gate and the executability respect---no intact internal
processing constitutes autonomy on this scale when the resource floor or the
action lever is gone (the gate and execution disjuncts of Theorem~T5).
Finally, the ethics of manipulation supplies the
operative concept of autonomy-undermining influence in technological settings:
\citet{sussner2019} define online manipulation as the covert exploitation of
decision-making weaknesses, distinguished from persuasion (which offers
reasons) and coercion (which threatens). This tripartite distinction is the
direct ancestor of our $\alpha$/$\beta$/$\gamma$ intervention classes.

\subsection{Cybernetic, autopoietic, and enactive accounts}\label{rw:bio}

% [MaM: keep; Syn: may compress by half]
Ashby's law of requisite variety established the information-theoretic core of
self-maintenance: a regulator must match the variety of the disturbances it
absorbs \citep{ashby1956}. Autopoiesis theory generalized this to living
systems, defining autonomy as the operatively closed self-production of a
system whose boundary is maintained by its own dynamics
\citep{maturana_varela1980}. \citet{di_paolo2005} added \emph{adaptivity}---the
capacity to modulate activity when self-maintenance is endangered---thereby
making precariousness, rather than robustness, the mark of the autonomous.
\citet{barandiaran2009} distil the tradition into three conditions for
agency: individuation, interactional asymmetry (the system's influence on its
milieu exceeds the converse), and normativity or sense-making. These accounts
supply the substrate-neutral backbone of our framework---resource reproduction
as a gating channel ($A_{via}$), openness as an axiom (AX2)---but remain
qualitative: no estimator, no scale, no aggregation rule.

\subsection{Psychological measurement of autonomy and its loss}\label{rw:psych}

Self-determination theory provides the most developed empirical construct of
autonomy. \citet{ryan_deci2000} model motivation along a continuum of
internalization from external regulation to integrated self-endorsement;
autonomy-supportive contexts promote internalization while controlling
ones---extrinsic rewards, threats, imposed deadlines---systematically
undermine it. Two findings are load-bearing for us. First, the continuum
vindicates treating endorsement as a matter of degree ($A_{goal}$). Second,
the ``paradox of compelled autonomy'': interventions whose explicit aim is to
produce autonomous behavior, deployed through controlling means, backfire; our
Theorem~T8 generalizes this regularity, and our $\sigma$-class
(scaffolding) captures the empirically supported alternative.

Learned helplessness supplies the paradigm case of execution-channel collapse.
\citet{maier_seligman2016} review fifty years of evidence, including the
neuroscientific reversal showing that passivity is the default and control is
learned: after sufficient experience of response--outcome independence,
organisms cease to exercise control \emph{even when it is restored}. Capability
remains; the action--outcome lever and the belief in it are gone---exactly the
zero semantics of $A_{exec}$ in Theorem~T5.

The measurement sciences of self-knowledge frame our reflection channel and
dual-track protocol. \citet{fleming_lau2014} separate metacognitive
\emph{sensitivity} from \emph{bias}, providing the estimator family for the
fidelity term of $A_{ref}$; \citet{moore_obhi2012} review intentional binding
and show that the felt sense of agency is measurable yet systematically
dissociable from objective causal contribution---the warrant for treating
self-report as calibration data and hallucination detector, never as evidence
(cf.\ also \citealp{lavazza2016}, on operationalizing free will as a
measurable multidimensional construct). From dynamical systems psychology and
ecology we adopt the methodology of early-warning signals for critical
transitions \citep{lade_gross2012}, the template for Propositions~C1--C2
(catastrophic collapse and hysteresis in autonomy loss).

\subsection{Information-theoretic and causal instruments}\label{rw:tools}

% [Syn: keep full passage; MaM: may compress]
Each of our channels binds to an established instrument. Empowerment---the
channel capacity from an agent's actions to its future sensory input---is the
canonical intrinsic measure of control \citep{klyubin2005}; its conceptual
foundations and estimation issues are reviewed in \citet{salge2013}, and
\citet{salge_polani2017} deploy it as a robot-safety norm (an agent that
maintains its own empowerment does not strip humans of control channels).
Because empowerment is a joint agent--environment property, it absorbs the
conversion-factor problem that defeats resource-based measures of execution
capacity. For directed influence, directed information
\citep{massey1990} and its causally conditioned extension to networks with
feedback \citep{kramer2003}, together with transfer entropy
\citep{schreiber2000}, provide the closed-loop machinery underlying
$A_{org}$; Massey's Theorem~1 ($I(X^N\!\to\!Y^N)\le I(X^N;Y^N)$ with equality
iff no feedback) is precisely the mediation-vs-bypass distinction our
Theorem~T2 formalizes.
The information-theoretic study of autonomy itself has a direct lineage that
our origination channel continues: \citet{bertschinger2008} define and
measure autonomy as informational closure---an agent is autonomous to the
extent that its sensory future is determined by its own state rather than by
the environment's---and \citet{krakauer2020} develop an information theory of
individuality in the same closure tradition. \citet{ay_polani2008} work out
the causal-network version of the same idea, analyzing autonomy through the
balance of internally determined versus externally driven information flows
in causal networks---structurally the distinction our origination ratio
\eqref{eq:aorg} makes quantitative with causally conditioned directed
information, and the missing link between the closure tradition and the
directed-information machinery. All three stay at the level of a
single quantity: without the goal/execution/epistemic/reflection
decomposition, without an intervention taxonomy, without a historical clause,
and without an aggregation theory connecting parts to a whole. On the
incentive side, power-seeking analysis shows that optimal policies in MDPs
tend toward states of high option value \citep{turner2021}---the
decision-theoretic shadow of our T4 vertex result (an attacker concentrates
budget on the highest-leverage channel) and of $\Aexec$'s empowerment
semantics. The two programs are complementary: power-seeking says what
trained optimizers tend to seek; we measure what an influence takes away.
Maximum-entropy inverse reinforcement learning recovers effective goals from
behavior \citep{ziebart2008}, yielding the $U_{eff}$ side of $A_{goal}$.
Estimation at scale is covered by variational empowerment
\citep{mohamed2015} and by $k$-NN and neural mutual-information estimators
\citep{kraskov2004,belghazi2018} (Section~\ref{sec:identifiability}).
Causal emergence---macro-level descriptions can carry more effective
information than micro-level ones \citep{hoel2013}---is the closest
information-theoretic cousin of our origination ratio: both ask where in the
causal stack the effective source of behavior lives. And Bayesian persuasion
\citep{kamenica2011} is the game-theoretic counterpart of our $\alpha$ class:
a sender who commits to honest signal structures steers a receiver through
the receiver's own intact channel---which is precisely why persuasion, as
opposed to channel surgery, is autonomy-compatible here
(Theorem~T2), and why its covert support-restriction analogue is not.
The intervention semantics of our five-class taxonomy is that of do-calculus
and independent causal mechanisms \citep{peters2017}. Belief-state POMDPs
\citep{kaelbling1998} model the epistemic constraints under which a
self-model's estimate of its own autonomy can drift from the truth
(Theorem~T3b). Change-point detection \citep{truong2020} supplies the
sequential-analysis backbone of the autonomy-margin theorem (T6): slow
sub-threshold drift is detectable if and only if verification budgets exceed
the drift rate, which is the content of the $\delta^*$ bound.

The free energy principle \citep{friston2010,friston2023} and its Markov
blanket formalism would ground self-maintenance from first principles. We
deliberately do not build on it: \citet{bruineberg2022} show that Markov
blanket structure is non-unique and interest-relative (the ``pearl blanket''
problem), so a blanket cannot by itself deliver the \emph{normative} self
required by autonomy ascriptions; we take boundaries as given (engineering
individuation) and reserve a blanket interface for future unification.
Integrated information theory \citep{oizumi2014,tegmark2016} is the closest
formal neighbor in spirit---axioms about a system's causal structure, pushed
forward into measures---but targets phenomenal consciousness; we cite it as
the methodological precedent for our axiom-to-measure strategy.

\subsection{Characterizing and governing AI agents}\label{rw:ai-agents}

The 2024--2026 literature converges, from independent directions, on
multidimensional gradations of AI agency. \citet{dung2024} proposes a
threshold conception on which a system's agency profile is jointly determined
by goal-directedness, autonomy, efficacy (direct impact on the world),
long-term planning, and acting for reasons; current LLMs fail the profile, but
architectural compositions \citep[e.g.][]{park2023} point forward. The survey
literature systematizes the underlying engineering paradigm---LLM cores with
planning, memory, tool-use, and action modules \citep{wang2023,xi2023}---whose
component-wise attack surfaces are catalogued by \citet{su2025}.
\citet{kasirzadeh2025} characterize AI agents along four dimensions---autonomy,
efficacy, goal complexity, generality---with gradations feeding
``agentic profiles'' for differentiated governance. \citet{feng2025} define
five levels of agent autonomy indexed by the user's role (operator,
collaborator, consultant, approver, observer), argue that autonomy level is a
design decision separable from capability and environment, and propose
autonomy certificates. The level-of-autonomy (LOA) tradition in human factors
predates these schemes by four decades and is the most direct prior art on
\emph{grading} autonomy: Sheridan--Verplank's ten-level scale of human--computer
control allocation \citep{sheridan1978} and the
Parasuraman--Sheridan--Wickens model of types and levels of automation
\citep{parasuraman2000} are adopted engineering practice, with empirically
calibrated usability consequences. Three differences locate our contribution
relative to LOA rather than against it: LOA grades \emph{task allocation}
(who does what) at design time, while we measure \emph{retained autonomy
under influence} at run time; LOA levels are ordinal labels without
estimators, while each of our channels carries one; and the LOA settings are
a special case of our delegation pair $(\lambda,q)$ of Theorem~T7, so the
tradition's gradations are projections of the channel space, not rivals of it
(Proposition~\ref{prop:c0}). Formal agency criteria
\citep{kenton2023} and empirical work on induced preference shifts in
recommender systems \citep{carroll2022} are the closest agent-side and
user-side formalizations. Governance work now separates capability, permission,
and autonomy level explicitly \citep{capability_permission2026}, while
\citet{xing2026} challenge the agent construct itself, finding identity and
self-regulation dimensions missing from engineering accounts. Practice-side
governance standards now mandate traceability and audit trails for deployed
models \citep{nist2023}, which our causal-logging necessity result
(Corollary~\ref{cor:logs-necessary}) makes precise for the autonomy-audit
case. Related
positions include the dynamical threshold account of \citet{vanlier2024} and
the delegation-side analyses of \citet{danaher2018} (AI assistants and the
achievement gap) and \citet{coeckelbergh2023} (AI threats to citizens'
epistemic agency as threats to autonomy).

Two absences remain diagnostic even after the LOA tradition is counted.
First, \emph{evaluation}: agent benchmarks
\citep{liu2023} and their critiques \citep{kapoor2024} measure capability and
cost. The 2024--2026 evaluation practice has sharpened the capability side
considerably---METR's time-horizon methodology measures how long a task an
agent can complete with a stated success rate \citep{kwa2025}, and the
capability-versus-propensity distinction (what an agent \emph{can} do versus
what it reliably \emph{does}) is now standard evaluation methodology---but
what it measures is capability under benign scaffolding, not the retention or
loss of autonomy under adversarial influence: those are different questions,
and the discordant regime of our own external-baseline check (manipulation
that collapses $A$ while slightly \emph{improving} functional outcome,
\S\ref{sec:e2checks}) is exactly where capability metrics and autonomy
audits dissociate. Second,
\emph{composition}: none of the frameworks specifies how its dimensions
aggregate, how they trade off, or when any of them hitting zero annihilates
the whole. Our Proposition~\ref{prop:c0} addresses both absences at once: each framework
is exhibited as the projection of our channel space onto the special case in
which the influence source $I$ is a benign principal---which is why none of
them can express adversarial scenarios (\S\ref{rw:safety}) while our $I$ ranges
over arbitrary interveners, including the user herself and third-party
attackers. A standing risk for any autonomy \emph{metric} is Goodhart
pressure---once audited, the measure becomes a target
\citep{manheim2019}; the framework's answer is structural rather than
rhetorical: gaming the audit is itself a surgery on the measurement channel
or the evidence stream (a $\gamma$/$\beta$ act against the auditor's graph),
which the registration discipline and the black-box degradation rules of
Section~\ref{sec:identifiability} are built to expose; residual
gaming-of-the-estimator risk is registered in Section~\ref{sec:limitations}.
On the phenomenological side, we follow
\citet{chalmers2023}, \citet{butlin2023}, and \citet{long2024} in neither
asserting nor excluding machine consciousness; our dual-track protocol is
designed to be agnostic, using subjective signals only for calibration and
divergence detection.

\subsection{AI safety as the study of autonomy loss}\label{rw:safety}

% [MaM: compress the case list; Syn: keep]
The safety literature has accumulated, case by case, exactly the failure modes
our taxonomy organizes; we map the load-bearing ones and compress the rest.
Reward tampering---agents seizing control of their own reward signal---is
analyzed in causal-influence-diagram terms that parallel our $\beta$-class
surgery semantics \citep{amodei2016,pan2022,everitt2021}. Goal
misgeneralization shows capability and goal coming apart
\citep{langosco2022,shah2022}, unified under proxy alignment by
\citet{ngo2022}. Mesa-optimization exhibits the manipulation case as
engineering reality: an inner optimizer can be selected for \emph{deceptive}
alignment, passing every behavioral endorsement test while its goals were
installed from outside \citep{hubinger2019}. Goal drift in deployed LM agents
\citep{arike2025} is the empirical counterpart of the slow
sub-threshold drift regime of Theorem~T6, and sycophancy
\citep{sharma2023} is a live instance of our hallucination index $\AHI>0$.
Prompt injection \citep{greshake2023} and self-propagating agent worms
\citep{cohen2024} show observation-channel inputs producing action-level
hijack---the $\alpha$-masquerading-as-$\beta$ case that Theorem~T2's mediation
criterion separates from persuasion (susceptibility is measured by
\citealp{andriushchenko2024}); dark patterns at web scale
\citep{nouwens2020} are the $\gamma$-class channel manipulation of human
agents. On the defensive side, corrigibility
\citep{soares2015} is the technical core of our delegation-quality parameter
$q$; safe interruptibility \citep{orseau2016} is the exit-option condition of
our $\delta$ class seen from the designer's side (registered as a property of
the agent's $\Aexec$ environment, not of its autonomy);
constitutional AI instantiates an explicit second-order review loop
\citep{bai2022}, and LM self-models are learnable yet corruptible
\citep{binder2024}, supporting both the feasibility
of $A_{ref}$ in silicon and its confinement to the calibration track. The
conceptual map of value alignment \citep{gabriel2020} frames the normative
pluralism we quarantine into the application layer (the value-pluralism risk
is registered in Section~\ref{sec:objections}, not denied).

What the safety corpus lacks is a common yardstick: each failure mode is
studied in its own vocabulary, so the field cannot say which attack costs an
agent more autonomy, cannot allocate defense across channels
(Theorem~T4), and cannot state a margin below which autonomy is provably
retained (Theorem~T6).

\subsection{Delegation, incentives, and adversarial allocation}\label{rw:delegation}

Three further literatures formalize pieces of the delegation and
attack--defense geometry that our theorems treat, and the connections deserve
explicit statement. \emph{Assistance games and cooperative inverse
reinforcement learning} \citep{hadfieldmenell2016} give the game-theoretic
formalization of aligned delegation: the principal's behavior is evidence
about the true objective, and the agent's optimal policy incorporates that
evidence. Our delegation pair $(\lambda,q)$ of Theorem~T7 measures the same
relationship from the audit side: $q=1$ delegation (transparency,
reversibility, and endorsability all maximal) is the assistance ideal, and its equivalence to honest
information (T7(iv)) is the audit-level counterpart of CIRL's incentive
compatibility, while a $\beta_G$ implant is precisely the failure mode in
which the ``principal's'' signal is adversarial rather than informative about
the true objective. \emph{Contract theory} \citep{laffont_martimort2002}
prices delegation through incentive-compatibility and monitoring constraints;
the ungoverned share $m=\lambda(1-q)$ is instead an information-theoretic
measure of what delegation costs in terms of \emph{whose normative structure
drives behavior}---non-welfarist, and complementary to the utility-transfer
accounting of the principal--agent tradition. \emph{Security resource
allocation} has its own canonical model: the Colonel~Blotto game
\citep{roberson2006}, in which two sides simultaneously allocate budgets
across battlefields and equilibria are mixed. T4's vertex concentration with
water-filling defense differs structurally: it is a one-sided comparative
static of attack effort against a fixed defense under log-linear damage, so
vertex concentration is a property of the linear objective rather than an
equilibrium phenomenon, and the defender's equalization is the min-max
response. The two models carve different slices of the attack--defense
problem---Blotto: symmetric simultaneous allocation; T4: asymmetric
sequential audit---and where they overlap they are consistent.

\subsection{Composite indicators and scale-consistent aggregation}\label{rw:composite}

% [Syn: keep full passage; MaM: may compress to three sentences]
Multidimensional measurement faces a standing methodological objection:
aggregation is arbitrary. The composite-indicator literature answers it. The
OECD/JRC handbook makes scale type discipline mandatory: the admissible
aggregation of sub-indicators is constrained by their measurement scales, and
for strictly positive ratio-scale components geometric aggregation is the
scale-consistent choice---a characterization due to the social-choice
analysis of \citet{ebert_welsch2004} and inherited by the handbook
\citep{oecd2008}. Our representation theorem's measurement-theoretic half
(route B) is the bounded-cube, order-invariance form of that classical
characterization; what we add is the minimal-trio analysis (zero annihilation
and biseparability as consequences), the harmonic-mean isolation of the
invariance axiom's role, and the audit embedding (gate, weight protocols,
falsification windows, sensitivity obligations)---not the product form
itself (Section~\ref{sec:representation}). The non-compensability debate
has named opponents and we cite them as such: the Mazziotta--Pareto index
penalizes unbalanced profiles while staying compensatory in form
\citep{mazziotta2016}, and benefit-of-the-doubt/DEA weighting lets the data
choose the weights \citep{cherchye2007}; both reject the zero-annihilation
stance we adopt, so non-compensability is a conscious commitment against an
articulated alternative, discharged through sensitivity obligations rather
than assumed---and within our axiom set it is carried by the invariance trio
(strict monotonicity, coordinatewise scale invariance, homogeneity), from
which zero annihilation (AX5) follows as a consequence
(Remark~\ref{rem:routeb-minimal}). Our Representation Theorem
derives the same conclusion from within the framework: continuity, strict
coordinatewise monotonicity, homogeneity of degree one, and coordinatewise
scale invariance of the induced
order (AX6b) characterize the weighted geometric mean uniquely---and an
explicit harmonic-mean counterexample shows that without the invariance axiom
the characterization fails. The two derivations---measurement-theoretic from
outside, axiomatic from inside---converging on the product form is our answer
to the arbitrariness objection, with robustness reported across
axiom-relaxation bands of the elasticity $\varsigma$ (the
$\varsigma\in[0.5,1)$ band relaxing AX6b and $\varsigma=2$ relaxing AX5, on
equal footing; Section~\ref{sec:experiments}).

The measurement-theoretic foundation of this scale-type discipline is the
Krantz--Luce--Suppes--Tversky theory of measurement
\citep{krantz1971}: admissible transformations are fixed by the scale type
of the underlying quantity, and an aggregation form is meaningful only if
invariant under them. The same normative burden that our AX5 carries has been
fought most explicitly in multidimensional poverty measurement, where
non-compensatory zero-annihilating aggregation is adopted as a conscious
axiomatic stance and its cost is standard material \citep{alkire2011}; we
adopt the stance with the same explicitness and discharge it through
sensitivity obligations rather than through silence.

\subsection{Summary of gaps}\label{rw:gaps}

Five structural gaps organize the field. (1) Philosophy has the concepts but no
measures: fifty years of analysis of endorsement, history, and relational
constitution stop short of estimators \citep{frankfurt1971,fischer_ravizza1998,mackenzie_stoljar2000}.
(2) AI governance has gradations but no theory: level schemes and dimension
profiles are calibrated neither to each other nor to any estimator, and all
presume a benign principal \citep{feng2025,kasirzadeh2025,dung2024}. (3) The
information-theoretic instruments are normatively inert: empowerment, directed
information, and IRL know nothing of endorsement, manipulation, or history
\citep{klyubin2005,massey1990,ziebart2008}. (4) AI safety has the failure
cases but no common scale on which to compare, prioritize, or bound them
(\S\ref{rw:safety}). (5) The three substrates---human, animal, artificial---are
studied in mutually untranslated vocabularies \citep{ryan_deci2000,barandiaran2009,kasirzadeh2025}.
The present framework's contribution against these five gaps is scoped, not
sweeping: it addresses (1), (3), and (4) by formalization---a substrate-neutral
definition of autonomy as the degree to which behavior is determined by the
agent's own normative structure $(G,B,M)$, decomposed into five channels behind
a viability gate, each
bound to an estimator, aggregated by the product form the stated axioms
characterize (under commitments that are defended, not derived), and
governed by theorems for retention (T6), loss (T5), mediation (T2),
manipulation (T3), and delegation (T7); it addresses (2) only by a
\emph{projection relation} (Proposition~\ref{prop:c0}, under definitional
readings whose faithfulness is registered as hermeneutic, not mathematical,
residual uncertainty); and it addresses (5) only by a common \emph{definition
schema}, not a common scale---cross-substrate comparability is protocol-bound
(\S\ref{sec:objections}(viii)). Every theorem carries an
explicit falsification window exercised in Section~\ref{sec:experiments}. The
substrate scope is stated precisely rather than maximally: the
\emph{artificial} substrate is exercised by the benchmark at the
implementation-consistency level; the
\emph{human} substrate is instantiated by protocol with literature
triangulation and a stated validation path (limited experiments and
large-scale online data); the \emph{animal} substrate is a conceptual
possibility within the L0--L1 layers, not a coverage claim of this paper.

% =====================================================================
% Note: the \citet/\citep above rely on natbib; the sn-jnl template is compatible by default.
% If the target journal requires numbered style, switch to \cite and replace spbasic with the sn-mathphys style.
% =====================================================================
